\section{Enlaces del Proyecto y Documentación}

\subsection{Repositorio de GitHub}
El código fuente de la aplicación, el esquema de la base de datos y la documentación compartida del proyecto están alojados en GitHub.

\begin{itemize}
    \item \textbf{Enlace al repositorio:} \url{https://github.com/DeltaCeleste/Doble-M}
    \item \textbf{Ubicación del código fuente:} El código backend principal (API y funciones del sistema) se encuentra en la ruta dentro del repositorio:
    \begin{verbatim}
~/Doble-M/app/www/
    \end{verbatim}
    \textit{Nota: Los controladores de la API se ubican en \texttt{www/api/} y la lógica de base de datos en \texttt{www/utils/functions.php}.}
\end{itemize}

\subsection{Seguimiento de Tareas (Jira)}
La planificación de las Historias de Usuario (HU), la asignación de subtareas de backend y el seguimiento del flujo de trabajo de la metodología Scrum se gestionan a través de Jira.

\begin{itemize}
    \item \textbf{Enlace al tablero Jira:} \url{https://go-team-yudm9r0y.atlassian.net/jira/software/projects/SCRUM/boards/1?filter=assignee+\%3D+712020\%3Ad02ff67f-2ee4-4aad-8765-56a3f4501dd4&groupBy=subtask&atlOrigin=eyJpIjoiZTRjNjYyNWNmMGU1NDZlNzlkNzY0ZTE4NjVlYTIzYTMiLCJwIjoiaiJ9}
    \item \textit{Nota sobre el acceso:} Es posible que se requiera realizar una invitación al usuario específico para acceder al tablero; en caso de no tener acceso, enviar un correo a \textbf{\texttt{deltac@correo.ugr.es}}.
\end{itemize}